% Part: first-order-logic
% Chapter: model-theory
% Section: theory-of-m

\documentclass[../../../include/open-logic-section]{subfiles}

\begin{document}

\olfileid[hi]{mod}{bas}{thm}

\section{किसी \printtoken{S}{structure} का सिद्धांत}

प्रत्येक !!{structure}~$\Struct{M}$ कुछ !!{sentence}s को सत्य और कुछ को
असत्य करती है। जिन सभी !!{sentence}s को वह सत्य करती है, उनके समुच्चय को
उसका \emph{सिद्धांत} कहते हैं। वह समुच्चय वास्तव में एक सिद्धांत है,
क्योंकि उससे निष्पन्न होने वाली प्रत्येक बात उसके सभी प्रतिरूपों में सत्य होनी
चाहिए, जिनमें~$\Struct{M}$ भी सम्मिलित है।

\begin{defn}
  कोई !!a{structure}~$\Struct M$ दी गई हो। $\Struct{M}$ का
  \emph{सिद्धांत}, उन !!{sentence}s का समुच्चय $\Theory{M}$ है जो
  $\Struct{M}$ में सत्य हैं, अर्थात् $\Theory{M} =
  \Setabs{!A}{\Sat{M}{!A}}$.
\end{defn}

हम ``सिद्धांत'' पद का अनौपचारिक उपयोग उन !!{sentence}s के समुच्चयों के लिए
भी करते हैं जिनकी कोई अभिप्रेत व्याख्या हो, चाहे वे निगमन के अधीन संवृत हों या
नहीं।

\begin{prop}
प्रत्येक $\Struct{M}$ के लिए $\Theory{M}$ पूर्ण है।
\end{prop}

\begin{proof}
किसी भी !!{sentence}~$!A$ के लिए या तो $\Sat{M}{!A}$ अथवा
$\Sat{M}{\lnot !A}$, इसलिए या तो $!A \in \Theory{M}$ अथवा
$\lnot !A \in \Theory{M}$।
\end{proof}

\begin{prop}\ollabel{prop:equiv}
  यदि $\Struct{N} \models !A$ प्रत्येक $!A \in \Theory{M}$ के लिए, तो
  $\Struct{M} \elemequiv \Struct{N}$।
\end{prop}

\begin{proof}
चूँकि $\Sat{N}{!A}$ सभी $!A \in \Theory{M}$ के लिए, $\Theory{M} \subseteq
\Theory{N}$। यदि $\Sat{N}{!A}$, तो $\Sat/{N}{\lnot !A}$, इसलिए $\lnot !A
\notin \Theory{M}$। चूँकि $\Theory{M}$ पूर्ण है, $!A \in
\Theory{M}$। अतः $\Theory{N} \subseteq \Theory{M}$, और इसलिए
$\Struct{M} \elemequiv \Struct{N}$।
\end{proof}

\begin{rem}\ollabel{remark:R}
  $\Struct{R} = \langle\Real, <\rangle$ पर विचार करें। यह वह
  !!{structure} है जिसका प्रांत वास्तविक संख्याओं का समुच्चय $\Real$ है,
  और जिसकी !!{language} में केवल एक द्वि-स्थानीय !!{predicate} है जिसकी
  व्याख्या वास्तविक संख्याओं पर $<$ संबंध के रूप में होती है। स्पष्टतः
  $\Struct{R}$ !!{nonenumerable} है; किंतु $\Theory{R}$ के स्पष्टतः संगत
  होने से, ल्योवेनहाइम--स्कोलेम (L\"owenheim--Skolem) प्रमेय के अनुसार उसका
  कोई !!a{enumerable} प्रतिरूप है, मान लें $\Struct{S}$।
  \olref{prop:equiv} से $\Struct{R}
  \equiv \Struct{S}$। फिर भी $\Struct{R}$ और $\Struct{S}$ समरूपी नहीं हैं,
  अतः इससे दिखता है कि \olref[iso]{thm:isom} का विलोम सामान्यतः विफल होता
  है।
\end{rem}

\end{document}
